\documentclass[a4paper,11pt]{article}
\usepackage[utf8]{inputenc}
\usepackage[T1]{fontenc}
\usepackage[french]{babel}
\usepackage{lmodern}
\usepackage{amsmath,amssymb,amsthm}
\usepackage{mathrsfs}
\usepackage{enumitem}
\usepackage{multicol}

\setlist[enumerate]{label=\textbf{\textcolor{Couleur8}{\arabic*.}}, left=1em}

% Si vous voulez aussi modifier les listes enumerate imbriquées :
\setlist[enumerate,2]{label=\textcolor{Couleur8}{\alph*)}}
\setlist[enumerate,3]{label=\textcolor{Couleur8}{\roman*)}}
\setlist[enumerate,4]{label=\textcolor{Couleur8}{\Alph*)}}
\usepackage{graphicx}
\usepackage{xcolor}
\usepackage{tcolorbox}
\usepackage{geometry}
\usepackage{fancyhdr}
\usepackage{titlesec}
\usepackage{cancel}
\usepackage{ifthen}
\usepackage{tikz}
\usetikzlibrary{arrows.meta}
\usepackage{tkz-tab}
\usepackage{eurosym}

% OPTION PROF/ÉLÈVE
% Décommentez UNE des deux lignes suivantes :
\newboolean{prof}\setboolean{prof}{true}   % Version professeur (avec corrections des devoirs)
%\newboolean{prof}\setboolean{prof}{false}  % Version élève (sans corrections des devoirs)

\definecolor{Couleur1}{HTML}{4A5D7A} %Th
\definecolor{Couleur2}{HTML}{A5B5D6} %Prop
\definecolor{Couleur3}{HTML}{A8C68A} %Méthode
\definecolor{Couleur6}{HTML}{8A9CC1} %Def
\definecolor{Couleur7}{HTML}{E6B459} %Rq Voc
\definecolor{Couleur8}{HTML}{D36740} %Titres
\definecolor{correction}{HTML}{D36740} % Couleur pour les corrections
\definecolor{coopmathscorr}{HTML}{F15929} % Couleur corrections coopmaths

% Configuration des marges
\geometry{
	top=2cm,
	bottom=2cm,
	inner=1.5cm,
	outer=1.5cm,
}

% Police
\renewcommand{\familydefault}{\sfdefault}

% Compteurs
\newcounter{seancenum}
\newcounter{seancenumD}
\setcounter{seancenum}{1}
\setcounter{seancenumD}{1}

% Configuration des en-têtes et pieds de page
\pagestyle{fancy}
\fancyhf{}
\ifthenelse{\boolean{prof}}{
    \fancyhead[L]{\textcolor{Couleur8}{\textbf{Corrections Automatismes - Classe de seconde - Version Professeur}}}
}{
    \fancyhead[L]{\textcolor{Couleur8}{\textbf{Corrections Automatismes - Séance 18 - Classe de seconde}}}
}
\fancyhead[R]{\textcolor{Couleur8}{\textbf{2025-2026}}}
\fancyfoot[L]{\textcolor{Couleur8}{Mme Bertail et Mme Pinto}}
\fancyfoot[R]{\textcolor{Couleur8}{\thepage}}
\renewcommand{\headrulewidth}{1pt}
\renewcommand{\headrule}{\hbox to\headwidth{\color{Couleur8}\leaders\hrule height \headrulewidth\hfill}}

% Environnements personnalisés
\newtcolorbox{seance}[1]{
    colback=Couleur6!10,
    colframe=Couleur1!65,
    fonttitle=\bfseries\large,
    title=Correction Séance \theseancenum #1,
    left=5mm,
    right=5mm,
    top=3mm,
    bottom=3mm,
    arc=0mm,
    after=\stepcounter{seancenum}
}

\newtcolorbox{seanceD}[1]{
    colback=Couleur6!5,
    colframe=Couleur6!45,
    fonttitle=\bfseries\large,
    title=Correction Devoirs - Séance \theseancenumD #1,
    left=5mm,
    right=5mm,
    top=3mm,
    bottom=3mm,
    arc=0mm,
    after=\stepcounter{seancenumD}
}

\begin{document}

\setcounter{seancenum}{18}
\newpage
\begin{seance}{} %18

\begin{enumerate}
\item On écrit les masses en écriture scientifique pour les comparer :\\
- Planète 1 : $891{,}33\times 10^{23} = 8{,}913\,3\times 10^{25}$ kg\\
- Planète 2 : $449{,}85\times 10^{23} = 4{,}498\,5\times 10^{25}$ kg\\
- Planète 3 : $106{,}959\,6\times 10^{24} = 1{,}069\,596\times 10^{26}$ kg\\
- Planète 4 : $945{,}99\times 10^{21} = 9{,}459\,9\times 10^{23}$ kg\\

On a donc : $1{,}069\,596\times 10^{26} > 8{,}913\,3\times 10^{25} > 4{,}498\,5\times 10^{25} > 9{,}459\,9\times 10^{23}$\\
Donc il s'agit de la {\bfseries \color{correction}Planète 3} qui a la masse la plus importante.\\
La bonne réponse est la réponse {\bfseries \color{correction}B}.

\item Le coefficient multiplicateur associé à une baisse de $27\,\%$ est $1-0{,}27=0{,}73$.\\
Le coefficient multiplicateur réciproque est donc $\dfrac{1}{0{,}73}$.\\
On en déduit que le taux réciproque est $\dfrac{1}{0{,}73}-1$ ou $\dfrac{0{,}27}{0{,}73}$.\\
Le taux à appliquer pour que cet article retrouve son prix initial est donné par ${\color{correction}\boldsymbol{\dfrac{1}{0{,}73}-1}}$.\\
La bonne réponse est la réponse {\bfseries \color{correction}B}.

\item  $(-5{,}78)^2< (-5{,}8)^2$ car $5,8 > 5,78$

\item On peut simplifier l'expression :\\
$\begin{aligned}
\dfrac{35x^{5}}{\dfrac{7}{x^2}}&=35x^{5} \times \dfrac{x^2}{7}\\
&=\dfrac{35x^{7}}{7}\\
&=5x^{7}.
\end{aligned}$\\
La bonne réponse est la réponse {\bfseries \color{correction}D}.

\item On a :\\
$\begin{aligned}
A&=\dfrac{286}{100}+\dfrac{29}{1\,000}\\
&=2{,}86+0{,}029\\
&={\color{correction}\boldsymbol{2{,}889}}
\end{aligned}$

\end{enumerate}

\end{seance}

\end{document}